Trois charges ponctuelles sont placées aux sommets d'un triangle équilatéral de \qty{\v{a}}{cm} de côté : $q_1 = \qty{\v{q1}}{\micro\coulomb}$, $q_2 = \qty{\v{q2}}{\micro\coulomb}$ et $q_3 = \qty{\v{q3}}{\micro\coulomb}$.

\begin{enumerate}[label=\alph*)]
    \item Déterminez l'énergie potentielle électrique totale du système des trois charges.
    \item Déterminez le potentiel électrique au centre du triangle (situé à $a/\sqrt{3}$ de chaque sommet).
    \item Quel travail un agent extérieur doit-il fournir pour amener une quatrième charge $q_4 = \qty{\v{q4}}{nC}$ de l'infini jusqu'au centre ?
\end{enumerate}
